\chapter{Lee--Yang 零点与热力学非解析性}
\label{chap:phys-sm-partition-function-zeros-phase-transitions}

有限体积零点与解析极限的关系说明：只有零点逼近物理参数轴，热力学极限才可能失去解析性。对铁磁 Ising 模型，还能进一步精确决定这些零点的几何位置。先在有限一般图上建立多变量 Lee--Yang 零点禁区，再把均匀外场零点限制到单位圆，并说明零点测度如何控制磁化和相变。这些结论不依赖二维方格的可解性；Onsager 自由能与 Yang 磁化来自另一条零场精确解路线。

\section{有限图铁磁 Ising 模型}
\label{sec:phys-sm-finite-zeros}

设 \(G=(\mathsf V,\mathsf E)\) 为有限无向图，自旋 \(\sigma_x\in\{\pm1\}\)，耦合 \(J_{xy}\ge0\)。Hamilton 量取
\begin{equation}
 H_G(\sigma)
 =-\sum_{\{x,y\}\in\mathsf E}J_{xy}\sigma_x\sigma_y
 -\sum_{x\in\mathsf V}h_x\sigma_x.
 \label{eq:phys-sm-general-ising-hamiltonian}
\end{equation}
定义 \(K_{xy}=\beta J_{xy}\) 以及局域复场变量
\begin{equation}
 z_x=\ee^{-2\beta h_x}.
 \label{eq:phys-sm-lee-yang-variable}
\end{equation}
把 \(\exp(\beta h_x)\) 等非零因子提出后，有限图配分函数成为每个 \(z_x\) 次数至多为一的多仿射多项式 \(\Phi_G(\{z_x\})\)。

单条铁磁边已经体现零点禁区的方向。忽略非零常数，二自旋边多项式可写成
\begin{equation}
 P_K(z_1,z_2)=1+z_1z_2+t(z_1+z_2),
 \qquad t=\ee^{-2K}\in(0,1].
 \label{eq:phys-sm-two-spin-edge-poly}
\end{equation}
若 \(P_K=0\)，则
\[
 z_2=-\frac{1+t z_1}{z_1+t}.
\]
直接计算
\begin{equation}
 |1+t z_1|^2-|z_1+t|^2
 =(1-t^2)(1-|z_1|^2).
 \label{eq:phys-sm-two-spin-stability}
\end{equation}
所以当 \(|z_1|<1\) 且 \(K>0\) 时必有 \(|z_2|>1\)。一条铁磁边不可能在双单位圆盘中产生零点。

把局部结论推广到任意有限图需要 Asano--Ruelle contraction lemma。该引理是独立的多变量稳定多项式定理：它说明在逐边合并变量的特定 contraction 操作下，上述单位圆盘零点禁区保持不变。简单代入 \(z_i=z_j\) 与 contraction 并不等价；从单边稳定性推广到一般图所需的保持定理正是 Asano--Ruelle lemma。

\begin{theorem}[多变量 Lee--Yang 定理]
\label{thm:phys-sm-multivariate-lee-yang}
对\eqnrefs{eq:phys-sm-general-ising-hamiltonian} 的任意有限铁磁 Ising 图，若所有 \(|z_x|<1\)，则
\[
 \Phi_G(\{z_x\})\neq0.
\]
由全自旋翻转对称性，若所有 \(|z_x|>1\)，配分函数同样不为零。
\end{theorem}

\begin{note}[独立数学输入的边界]
\eqnrefs{eq:phys-sm-two-spin-stability} 可由二自旋多项式直接计算；从单边稳定性到任意图的逐边 contraction 则依赖 Asano--Ruelle lemma。两者的证明层级必须分开。
\end{note}

\begin{derivation}[Asano--Ruelle contraction 引理与多变量 Lee--Yang 定理]
\emph{策略。}Asano--Ruelle contraction 引理把单边稳定性\eqnrefs{eq:phys-sm-two-spin-stability} 逐边推广到任意有限图。核心操作是"contract"：把双变量多项式中 \(z_1,z_2\) 同时替换为 \(z\)，在特定条件下保持零点禁区。

\emph{第一步：Asano contraction 的代数表述。}设 \(P(z_1,z_2)\) 是双变量多项式，在 \(\mathbb D^2=\{|z_1|<1,|z_2|<1\}\) 内无零点（即 \(P\) 在双单位圆盘上稳定）。定义 contract
\begin{equation*}
 (CP)(z)=P(z,z).
\end{equation*}
\emph{引理（Asano--Ruelle）：}\(CP\) 在 \(\mathbb D=\{|z|<1\}\) 内无零点。

\emph{证明。}反证法。设存在 \(z_0\in\mathbb D\) 使 \(P(z_0,z_0)=0\)。把 \(P\) 按 \(z_2\) 展开：
\begin{equation*}
 P(z_1,z_2)=\sum_{k=0}^n a_k(z_1)\,z_2^k.
\end{equation*}
对固定 \(z_1\in\mathbb D\)，\(P(z_1,\cdot)\) 在 \(\mathbb D\) 内无零点（假设），故所有根 \(w_j(z_1)\) 满足 \(|w_j|\ge1\)。由根与系数关系，\(P(z_1,z_2)=a_n(z_1)\prod_j(z_2-w_j(z_1))\)。

关键步骤：对每个 \(z_1\in\mathbb D\)，定义 Rouché 变换
\begin{equation*}
 \tilde P(z_1,\zeta)=\zeta^n P(z_1,\zeta^{-1})=a_0(z_1)+a_1(z_1)\zeta+\cdots+a_n(z_1)\zeta^n.
\end{equation*}
\(\tilde P\) 的根为 \(\zeta_j=1/w_j\)，故 \(|\zeta_j|\le1\)。由假设 \(P\) 在 \(\mathbb D^2\) 无零点，当 \(|z_1|<1\) 时 \(P(z_1,z_2)\ne0\) 对 \(|z_2|<1\)，即 \(w_j\notin\mathbb D\)，故 \(\zeta_j\notin\overline{\mathbb D}^c\)，即 \(|\zeta_j|\le1\)。

现在 \(P(z_0,z_0)=0\) 意味着 \(\tilde P(z_0,1/z_0)=z_0^{-n}P(z_0,z_0)=0\)。但 \(|1/z_0|>1\)，而 \(\tilde P(z_0,\cdot)\) 的根 \(\zeta_j\) 满足 \(|\zeta_j|\le1\)，故 \(1/z_0\) 不是根——矛盾。

\emph{第二步：Ruelle 推广到多变量。}Ruelle 把 Asano contraction 推广到多变量：设 \(P(z_1,\ldots,z_n)\) 在 \(\mathbb D^n\) 内无零点。对任意 \(i\ne j\)，contract \(z_i,z_j\to z\)（即令 \(z_i=z_j=z\)）后，\(P(z,\ldots,z,\ldots,z,\ldots)\) 在 \(\mathbb D\) 内无零点。证明对 \(z_i\) 固定其余变量，应用第一步的双变量结果。

\emph{第三步：从单边到一般图。}对有限铁磁图 \(G=(\mathsf V,\mathsf E)\)，配分函数 \(\Phi_G(\{z_x\})\) 可逐边构造：
\begin{itemize}
 \item 初始：每个顶点 \(x\) 贡献因子 \((1+z_x)\)（外场部分），在 \(\mathbb D^{|\mathsf V|}\) 内无零点。
 \item 逐边添加：对边 \(\{x,y\}\) 添加铁磁耦合因子 \(P_{K_{xy}}(z_x,z_y)=1+z_xz_y+t_{xy}(z_x+z_y)\)（\eqnrefs{eq:phys-sm-two-spin-edge-poly}）。由\eqnrefs{eq:phys-sm-two-spin-stability}，\(P_{K_{xy}}\) 在 \(\mathbb D^2\) 内无零点。
 \item 合并：添加边 \(\{x,y\}\) 后，新多项式是旧多项式与 \(P_{K_{xy}}\) 的乘积。乘积在 \(\mathbb D^n\) 内无零点当且仅当每个因子无零点。但合并后 \(z_x,z_y\) 成为同一多项式的变量——需用 Asano--Ruelle contraction 保证合并后仍稳定。
\end{itemize}
具体地，逐边构造中每步产生的新变量（如合并两个独立子图时共享的顶点变量）由 Ruelle 推广保证稳定性传递。最终 \(\Phi_G\) 在 \(\mathbb D^{|\mathsf V|}\) 内无零点，即多变量 Lee--Yang 定理。

\emph{第四步：圆定理。}取均匀场 \(z_x=z\)，\(\Phi_G(z,\ldots,z)=\Phi_G(z)\) 是单变量多项式。由多变量定理，\(\Phi_G(z)\ne0\) 当 \(|z|<1\)；由自旋翻转对称 \(\Phi_G(z)=z^{|\mathsf V|}\Phi_G(1/z)\)，\(\Phi_G(z)\ne0\) 当 \(|z|>1\)。故零点全在 \(|z|=1\)，即 Lee--Yang 圆定理。

\emph{物理意义。}Asano--Ruelle 引理的物理核心是"铁磁耦合保持零点禁区"：单条铁磁边把零点排除在双单位圆盘外，逐边合并时 contraction 操作保持这一禁区不变。铁磁性（\(J_{xy}\ge0\)）是关键——反铁磁耦合（\(J_{xy}<0\)）使 \(t=\ee^{-2K}>1\)，\eqnrefs{eq:phys-sm-two-spin-stability} 中 \((1-t^2)\) 变号，禁区不再保持。这解释了为何 Lee--Yang 圆定理只适用于铁磁 Ising 模型。Asano--Ruelle 引理在多变量稳定多项式理论中是基本工具，还应用于量子自旋玻璃零点分布和连续变量 Lee--Yang 定理（Lieb--Seiringer 定理）。零点分布的禁区可总结为
\begin{equation*}
\{z\in\CC:\Phi_G(z)\ne0\}\supseteq\{z:|z|<1\}\cup\{z:|z|>1\},\quad\text{故零点}\subseteq\{|z|=1\}.
\end{equation*}



\end{derivation}

\section{均匀场与 Lee--Yang 圆定理}
\label{sec:phys-sm-circle-theorem}

取均匀场 \(h_x=h\)，于是所有 \(z_x=z=\ee^{-2\beta h}\)。多变量定理立即给出：当 \(|z|<1\) 或 \(|z|>1\) 时配分函数无零点。因此所有有限体积零点只能位于
\begin{equation}
 |z|=1.
 \label{eq:phys-sm-lee-yang-circle-general}
\end{equation}
这就是 Lee--Yang 圆定理。物理实场 \(h\in\mathbb R\) 对应正实 \(z>0\)，而 \(h=0\) 对应 \(z=1\)。因此对铁磁 Ising 模型，有限体积零点逼近物理轴的唯一可能位置就是单位圆与正实轴的交点 \(z=1\)。

把零点写成 \(z_\alpha=\ee^{\ii\theta_\alpha}\)，单位体积零点角密度若在热力学极限存在，则记为 \(\rho_{\rm LY}(\theta)\)。由零点 Cauchy 变换\eqnrefs{eq:phys-sm-zero-cauchy-transform}，均匀磁化可写成零点分布的边界值。若 \(N=|\mathsf V|\)，定义每自旋磁化
\begin{equation*}
 \mathfrak m_N=\frac{1}{N}\left\langle\sum_x\sigma_x\right\rangle
 =\frac{1}{N\beta}\frac{\partial\log Z}{\partial h}.
\end{equation*}
由于 \(z=\ee^{-2\beta h}\)，有
\begin{equation}
 \mathfrak m_N(z)=1-\frac{2z}{N}\frac{\partial\log P_N(z)}{\partial z}
 =1-2z\int_{-\pi}^{\pi}
 \frac{\rho_{{\rm LY},N}(\theta)}{z-\ee^{\ii\theta}}\,\dd\theta,
 \label{eq:phys-sm-magnetization-zero-transform}
\end{equation}
其中 \(P_N\) 是提出非零场因子后的零点多项式，\(\rho_{{\rm LY},N}(\theta)=N^{-1}\sum_\alpha\delta(\theta-\theta_\alpha)\) 是有限体积归一化角密度。

\begin{derivation}[Sokhotski--Plemelj 公式与 Cauchy 边界值]
\emph{策略。}Sokhotski--Plemelj 公式给出 Cauchy 型积分在积分路径上的边界值，是从零点密度到磁化跳跃的关键工具。

\emph{Cauchy 型积分。}设 \(\varphi(t)\) 在实轴上 Hölder 连续（\(|\varphi(t)-\varphi(s)|\le C|t-s|^\alpha\)，\(0<\alpha\le1\)）。定义 Cauchy 型积分
\begin{equation*}
 \Phi(z)=\frac{1}{2\pi i}\int_{-\infty}^{\infty}\frac{\varphi(t)}{t-z}\,\dd t,\qquad z\notin\mathbb R.
\end{equation*}
\(\Phi\) 在上下半平面分别全纯。问题：当 \(z\) 从上下半平面趋于实轴上一点 \(x_0\) 时，\(\Phi\) 的边界值是什么？

\emph{Sokhotski--Plemelj 公式。}
\begin{equation}
 \Phi^\pm(x_0)=\lim_{\varepsilon\downarrow0}\Phi(x_0\pm i\varepsilon)
 =\pm\frac12\varphi(x_0)+\frac{1}{2\pi i}\,\mathrm{p.v.}\!\int_{-\infty}^{\infty}\frac{\varphi(t)}{t-x_0}\,\dd t,
 \label{eq:phys-sm-sokhotski-plemelj}
\end{equation}
其中 \(\mathrm{p.v.}\) 为 Cauchy 主值积分。两式相减给出跳跃公式
\begin{equation}
 \Phi^+(x_0)-\Phi^-(x_0)=\varphi(x_0).
 \label{eq:phys-sm-plemelj-jump}
\end{equation}

\emph{证明。}把 \(\Phi(x_0+i\varepsilon)\) 拆成三部分：在 \(|t-x_0|<\delta\) 的小邻域内、邻域外、以及一个补偿项。
\begin{equation*}
 \Phi(x_0+i\varepsilon)=\frac{1}{2\pi i}\int_{|t-x_0|<\delta}\frac{\varphi(t)}{t-x_0-i\varepsilon}\,\dd t+\frac{1}{2\pi i}\int_{|t-x_0|\ge\delta}\frac{\varphi(t)}{t-x_0-i\varepsilon}\,\dd t.
\end{equation*}
第二项在 \(\varepsilon\to0\) 时直接趋于 \(\frac{1}{2\pi i}\int_{|t-x_0|\ge\delta}\frac{\varphi(t)}{t-x_0}\dd t\)。第一项改写为
\begin{align*}
 \frac{1}{2\pi i}\int_{|t-x_0|<\delta}\frac{\varphi(t)-\varphi(x_0)}{t-x_0-i\varepsilon}\,\dd t
 +\frac{\varphi(x_0)}{2\pi i}\int_{|t-x_0|<\delta}\frac{\dd t}{t-x_0-i\varepsilon}.
\end{align*}
由 Hölder 条件，第一部分被积函数有界 \(C|t-x_0|^\alpha/|t-x_0-i\varepsilon|\le C|t-x_0|^{\alpha-1}\)，可积（\(\alpha>0\)），\(\varepsilon\to0\) 时趋于 \(\frac{1}{2\pi i}\int_{|t-x_0|<\delta}\frac{\varphi(t)-\varphi(x_0)}{t-x_0}\dd t\)。第二部分直接计算：
\begin{equation*}
 \int_{x_0-\delta}^{x_0+\delta}\frac{\dd t}{t-x_0-i\varepsilon}
 =\log(t-x_0-i\varepsilon)\Big|_{x_0-\delta}^{x_0+\delta}
 =\log(\delta-i\varepsilon)-\log(-\delta-i\varepsilon).
\end{equation*}
当 \(\varepsilon\downarrow0\) 时，\(\log(\delta-i\varepsilon)\to\log\delta\)，\(\log(-\delta-i\varepsilon)\to\log\delta-i\pi\)（从下半平面绕过支点），故积分趋于 \(i\pi\)。因此
\begin{equation*}
 \Phi^+(x_0)=\frac{\varphi(x_0)}{2}+\frac{1}{2\pi i}\,\mathrm{p.v.}\!\int\frac{\varphi(t)}{t-x_0}\,\dd t.
\end{equation*}
类似地 \(\Phi^-(x_0)\) 从 \(\varepsilon<0\) 方向取极限，\(\log(-\delta+i\varepsilon)\to\log\delta+i\pi\)，积分趋于 \(-i\pi\)，给出 \(\Phi^-(x_0)=-\frac{\varphi(x_0)}{2}+\frac{1}{2\pi i}\mathrm{p.v.}\int\frac{\varphi(t)}{t-x_0}\dd t\)。

\emph{物理意义。}Sokhotski--Plemelj 公式的物理核心是"Cauchy 核 \(1/(t-z)\) 在实轴上有 \(\pm i\pi\delta(t)\) 的分布极限"：
\begin{equation*}
 \lim_{\varepsilon\downarrow0}\frac{1}{t-x_0\pm i\varepsilon}=\mathrm{p.v.}\frac{1}{t-x_0}\mp i\pi\delta(t-x_0).
\end{equation*}
这一分布恒等式在量子场论中给出推迟/超前 Green 函数的谱表示（Kramers--Kronig 关系），在统计力学中给出零点密度与磁化跳跃的精确联系\eqnrefs{eq:phys-sm-mag-jump}。跳跃公式\eqnrefs{eq:phys-sm-plemelj-jump} 表明 Cauchy 型积分的"不连续跳跃"完全由密度函数 \(\varphi\) 决定——这是 Riemann--Hilbert 问题的基本工具，把"从边界值反解全纯函数"化为积分方程。
\end{derivation}

\begin{derivation}[零点密度与磁化跳跃]
设热力学极限零点密度 \(\rho_{\rm LY}(\theta)\) 在 \(\theta=0\) 邻域有边界值，并把\eqnrefs{eq:phys-sm-magnetization-zero-transform}中的 Cauchy 积分记为
\[
 F(z)=\int_{-\pi}^{\pi}
 \frac{\rho_{\rm LY}(\theta)}{z-\ee^{\ii\theta}}\,\dd\theta,
\]
于是热力学极限磁化可写作 \(\mathfrak m(z)=1-2zF(z)\)。当 \(z\) 从单位圆内外趋于 \(\ee^{\ii\varphi}\) 时，把核分式改写为
\[
 \frac1{z-\ee^{\ii\theta}}
 =\frac{\ee^{-\ii\theta}}{z\ee^{-\ii\theta}-1},
\]
并使用 Sokhotski--Plemelj 公式
\[
 \lim_{\varepsilon\downarrow0}
 \int_{-\pi}^{\pi}
 \frac{\rho_{\rm LY}(\theta)}{\theta-\varphi\pm\ii\varepsilon}\,\dd\theta
 =\mathrm{p.v.}\!\int_{-\pi}^{\pi}
 \frac{\rho_{\rm LY}(\theta)}{\theta-\varphi}\,\dd\theta
 \mp\ii\pi\rho_{\rm LY}(\varphi),
\]
其中主值部分在内外两侧相同，虚部贡献相差一个符号。由此 Cauchy 边界值之差为
\[
 F_{\rm in}(\varphi)-F_{\rm out}(\varphi)
 =-2\pi\ee^{-\ii\varphi}\rho_{\rm LY}(\varphi).
\]
代入 \(\mathfrak m(z)=1-2zF(z)\)，常数项 \(1\) 与因子 \(2z\) 在两侧取相同边界值 \(2\ee^{\ii\varphi}\)，因此相减后只剩
\begin{equation*}
 \mathfrak m_{\rm in}(\varphi)-\mathfrak m_{\rm out}(\varphi)
 =-2\ee^{\ii\varphi}
 \bigl[F_{\rm in}(\varphi)-F_{\rm out}(\varphi)\bigr]
 =4\pi\rho_{\rm LY}(\varphi).
\end{equation*}
亦即
\begin{equation}
 \mathfrak m_{\rm in}(\varphi)-\mathfrak m_{\rm out}(\varphi)
 =4\pi \rho_{\rm LY}(\varphi).
 \label{eq:phys-sm-mag-jump}
\end{equation}
现在取 \(\varphi=0\)，于是 \(z=\ee^{\ii\varphi}=1\)。由 \(z=\ee^{-2\beta h}\)，实场 \(h>0\) 给出 \(0<z<1\)，对应单位圆内侧；实场 \(h<0\) 给出 \(z>1\)，对应单位圆外侧。因此内侧边界 \(h\downarrow0^+\)，外侧边界 \(h\uparrow0^-\)，代入得
\[
 \mathfrak m(0^+)-\mathfrak m(0^-)=4\pi\rho_{\rm LY}(0).
\]
若自旋翻转对称相在低温分裂为 \(\pm\mathfrak m_{\rm sp}\)，则 \(\mathfrak m(0^+)=\mathfrak m_{\rm sp}\)、\(\mathfrak m(0^-)=-\mathfrak m_{\rm sp}\)，左端为 \(2\mathfrak m_{\rm sp}\)，因此
\[
 \rho_{\rm LY}(0)=\frac{\mathfrak m_{\rm sp}}{2\pi}.
\]

\textbf{物理意义。}Lee--Yang 零点密度 $\rho_{\rm LY}(h)$ 把配分函数的零点分布与磁化曲线直接联系：$\rho_{\rm LY}(0)$ 正比于自发磁化跳跃。在相变点 $h=0$，零点密度从零跳到 $\mathfrak m_{\rm sp}/(2\pi)$——这是 Yang--Lee 边缘定理的实轴版本。实例：二维 Ising 的 $\mathfrak m_{\rm sp}=(1-\sinh^{-4}2K)^{1/8}$，在 $T_c$ 处 $\rho_{\rm LY}(0)$ 从零连续增长，对应二阶相变；一维 Ising 无自发磁化，$\rho_{\rm LY}(0)=0$。
\end{derivation}

\eqnrefs{eq:phys-sm-mag-jump} 给出一个清楚的物理分类。若 \(\rho_{\rm LY}(0)>0\)，零场磁化发生跳跃；若零点只逼近 \(z=1\) 但密度在那里趋零，则可能得到连续相变；若零点与 \(z=1\) 保持有限角隙，则零场附近保持解析。

\section{一维链作为零点禁区的可解例子}
\label{sec:phys-sm-1d-ising}

一维周期 Ising 链提供最小的非平凡检验。取
\[
 H_N=-J\sum_{j=1}^{N}\sigma_j\sigma_{j+1}
 -h\sum_{j=1}^{N}\sigma_j,
 \qquad \sigma_{N+1}=\sigma_1,
\]
转移矩阵
\[
 \mathsf T_{\rm 1D}=
 \begin{pmatrix}
 \ee^{K+\beta h}&\ee^{-K}\\
 \ee^{-K}&\ee^{K-\beta h}
 \end{pmatrix},
 \qquad K=\beta J.
\]
其本征值记为 \(\Lambda_\pm^{\rm 1D}\)，并有
\[
 \Lambda_\pm^{\rm 1D}
 =\ee^K\left(
 \cosh\beta h
 \pm\sqrt{\sinh^2\beta h+\ee^{-4K}}
 \right),
\]
且 \(Z_N=(\Lambda_+^{\rm 1D})^N+(\Lambda_-^{\rm 1D})^N\)。令 \(h=\ii\varphi/(2\beta)\)，即 \(z=\ee^{-\ii\varphi}\) 位于单位圆。零点要求 \((\Lambda_+^{\rm 1D}/\Lambda_-^{\rm 1D})^N=-1\)，这只有在
\begin{equation}
 \left|\sin\frac{\varphi}{2}\right|
 \ge \ee^{-2K}
 \label{eq:phys-sm-1d-zero-gap-condition}
\end{equation}
时才可能满足。因此单位圆在 \(z=1\) 附近存在角隙
\begin{equation}
 \varphi_0(K)=2\arcsin(\ee^{-2K})>0
 \qquad (T>0).
 \label{eq:phys-sm-1d-gap-angle}
\end{equation}
\eqnrefs{eq:phys-sm-1d-gap-angle} 严格排除的是零场附近沿复外场方向的 Lee--Yang 奇性；它本身还不能单独排除温度方向的非解析性。后者可直接由零场最大本征值得到。取 \(h=0\) 时
\[
 \Lambda_+^{\rm 1D}=2\cosh K,
 \qquad
 \Lambda_-^{\rm 1D}=2\sinh K,
\]
所以对任意有限 \(K\)
\[
 \left|\frac{\Lambda_-^{\rm 1D}}{\Lambda_+^{\rm 1D}}\right|
 =\tanh K<1.
\]
因此热力学极限自由能严格由唯一主导本征值给出
\begin{equation}
 -\beta f_{\rm 1D}=\log(2\cosh K),
 \label{eq:phys-sm-1d-free-energy}
\end{equation}
它对所有 \(0<T<\infty\) 都解析。Lee--Yang 角隙与这一独立谱论证合在一起，才完整说明一维最近邻 Ising 链没有有限温相变。


\begin{derivation}[一维 Lee--Yang 零点的相位量子化]
令
\[
 u=\sin\frac{\varphi}{2},
 \qquad
 a=\ee^{-2K}.
\]
在单位圆上，\(\sinh(\beta h)=\ii u\)，故当 \(|u|\ge a\) 时根号为纯虚数。定义
\[
 r=\sqrt{u^2-a^2},
 \qquad
 c=\cos\frac{\varphi}{2},
\]
则
\[
 \Lambda_\pm^{\rm 1D}=\ee^K(c\pm\ii r).
\]
两者模长相同，因此可写
\[
 \frac{\Lambda_+^{\rm 1D}}{\Lambda_-^{\rm 1D}}
 =\ee^{2\ii\psi(\varphi)},
 \qquad
 \tan\psi(\varphi)=\frac{r}{c}.
\]
有限链零点条件成为
\[
 \ee^{2\ii N\psi}=-1,
\]
从而
\[
 \psi(\varphi_j)=\frac{(2j+1)\pi}{2N},
 \qquad j\in\mathbb Z.
\]
这给出有限 \(N\) 零点在允许弧上的离散相位量子化；而 \(|u|=a\) 正是允许弧的端点，因此直接恢复\eqnrefs{eq:phys-sm-1d-zero-gap-condition}--\eqnrefs{eq:phys-sm-1d-gap-angle}。

\emph{零点密度的热力学极限。}相邻零点相位差
\begin{equation*}
 \Delta\varphi
 =\varphi_{j+1}-\varphi_j
 \approx\frac{\pi}{N\,\psi'(\varphi_j)},
\end{equation*}
故单位体积零点角密度
\begin{equation*}
 \rho_{\rm LY}^{(N)}(\varphi)
 =\frac{1}{N\,\Delta\varphi}
 =\frac{\psi'(\varphi)}{\pi}
 =\frac{1}{\pi}\frac{\dd}{\dd\varphi}\arctan\!\frac{r}{c}.
\end{equation*}
显式求导并化简 \(r^2+c^2=u^2-a^2+c^2=1-a^2\)，得
\begin{equation*}
 \rho_{\rm LY}(\varphi)
 =\frac{1}{2\pi}\frac{|u|\,|u'|}{1-a^2}
 =\frac{1}{2\pi}\frac{|\sin(\varphi/2)\cos(\varphi/2)|}{1-\ee^{-4K}}
 =\frac{|\sin\varphi|}{4\pi(1-\ee^{-4K})},
\end{equation*}
仅在允许弧 \(\varphi_0\le|\varphi|\le\pi-\varphi_0\) 上非零，端点处密度按平方根发散（van Hove 奇性）。

\emph{与二维临界行为的对比。}一维零点密度在 \(\varphi=0\) 处严格为零（角隙 \(\varphi_0>0\)），由\eqnrefs{eq:phys-sm-mag-jump} 给出 \(\mathfrak m_{\rm sp}=0\)，无自发磁化。二维方格在 \(T<T_c\) 时角隙闭合，\(\rho_{\rm LY}(0)>0\)，自发磁化非零；\(T=T_c\) 时角隙刚刚闭合且 \(\rho_{\rm LY}(\varphi)\sim|\varphi|^{\delta-1}\)（\(\delta=15\) 为磁化临界指数），对应连续相变。Lee--Yang 圆定理本身只固定零点位于单位圆，具体密度由模型维度与谱结构决定：一维可由上述相位量子化显式算出，二维则需要 Onsager 谱或 Kac--Ward 行列式才能得到 \(\rho_{\rm LY}\) 的精确形式。

\emph{Fisher 零点对照。}若改以温度为复参数（Fisher 零点），一维链零点由 \(\Lambda_+^{\rm 1D}=\pm\ii\Lambda_-^{\rm 1D}\) 给出，位于 \(z=\ee^{-2K}\) 复平面虚轴，热力学极限下密度 \(\rho_{\rm F}(z)\) 在虚轴上均匀分布，物理实轴 \(K>0\) 上无 Fisher 零点积聚，与一维无相变一致。二维方格的 Fisher 零点则在 \(K=K_c\) 处积聚于实轴，与 Lee--Yang 零点在 \(z=1\) 处的积聚互补地刻画同一相变：前者描述温度方向奇性，后者描述外场方向奇性。

\emph{物理意义。}Lee--Yang 零点的相位量子化把"配分函数的零点在哪里"这一复分析问题转化为"允许弧上离散相位如何分布"的几何问题。物理上，零点密度 \(\rho_{\rm LY}(\varphi)\) 直接控制自发磁化：由\eqnrefs{eq:phys-sm-mag-jump}，\(\mathfrak m_{\rm sp}=\int_0^\pi\rho_{\rm LY}(\varphi)\dd\varphi\)，一维角隙 \(\varphi_0>0\) 使 \(\rho_{\rm LY}(0)=0\)，故无自发磁化；二维临界温度以下角隙闭合，\(\rho_{\rm LY}(0)>0\)，自发磁化非零。零点在允许弧端点的平方根发散（van Hove 奇性）是态密度的普适边缘行为，在光子晶体、声子谱和能带边缘中反复出现。相位量子化还提供了有限体积相变的"预兆"：虽然有限系统永远解析，零点随 \(N\) 增大向物理轴的逼近速率和密度分布已编码了热力学极限的相变性质——这是有限大小标度理论的基础。
\end{derivation}

\begin{exbox}[一维零点与二维相变的逻辑差别]
一维结论并不是“转移矩阵只有两个本征值，所以不能相变”。真正决定解析性的是热力学极限中最大本征值之间是否能在物理参数上形成非解析竞争。有限温一维链的 Lee--Yang 零点始终保留\eqnrefs{eq:phys-sm-1d-gap-angle} 的有限角隙；二维方格在临界温度以下则允许零点压到 \(z=1\)。后者需要后续二维几何与精确谱结构才能定量算出。
\end{exbox}

\section{Lee--Yang 理论的适用域}

圆定理依赖铁磁耦合和 Ising 型二值自旋，不是任意统计模型的普遍零点定理；零点积聚导致非解析性的逻辑则更一般。在下面的零场主线中，Lee--Yang 理论只负责外场方向的解析结构。零场二维 Ising 的自由能将通过另一条路线求解：先把配分函数改写成偶子图，再利用平面对偶Duality、transfer matrix 与费米子化。两条路线最终在 Yang 自发磁化处汇合，因为 Yang 自发磁化既决定长距离关联，也通过\eqnrefs{eq:phys-sm-mag-jump} 决定 \(z=1\) 处的零点密度。

\emph{[圆定理的一般化：多变量与连续自旋]}
Lee--Yang 圆定理可沿两个方向推广：多变量耦合（Asano--Ruelle 多变量圆定理）和连续自旋（Simon--Lieb 推广）。分四步建立。

\emph{第一步：多变量 Lee--Yang。}对图 \(G=(V,E)\) 上 Ising 模型，每顶点有独立外场 \(h_i\)，对应变量 \(z_i=\ee^{-2\beta h_i}\)。多变量配分函数 \(Z(z_1,\ldots,z_N)\) 是多变量多项式。Asano--Ruelle 定理：若所有耦合 \(J_{ij}\ge0\)，则 \(Z\) 的零点满足"多变量圆条件"——对任意 \(i\)，固定其余 \(z_j\)（\(|z_j|<1\)），\(z_i\) 的零点满足 \(|z_i|>1\)。证明：逐顶点应用 Asano--Ruelle 收缩，每次把两变量多项式收缩为单变量，零点保持在适当圆外。多变量圆定理蕴含单变量圆定理（取所有 \(z_i=z\)）。

\emph{第二步：Simon--Lieb 连续自旋推广。}Simon--Lieb（1974）把圆定理推广到连续自旋分布。设自旋 \(\sigma_i\in\mathbb R\) 服从分布 \(d\mu(\sigma)=g(\sigma)\dd\sigma\)，配分函数 \(Z=\int\prod_i d\mu(\sigma_i)\ee^{\sum_{ij}J_{ij}\sigma_i\sigma_j+\sum_i h_i\sigma_i}\)。若 \(g\) 满足 Lee--Yang 性质——其 Fourier 变换 \(\hat g(k)\) 的零点全为实数——则 \(Z\) 的零点在 \(h_i\) 复平面上的位置+位于 \(\operatorname{Re} h_i=0\) 附近。Gaussian \(g(\sigma)\propto\ee^{-\sigma^2/2}\) 满足此条件（\(\hat g\propto\ee^{-k^2/2}\) 无零点），双峰 \(g\propto\delta(\sigma-1)+\delta(\sigma+1)\) 也满足（\(\hat g=\cos k\)，零点全实）。这把 Ising 的离散自旋推广到连续 Gaussian（\(O(1)\) 模型）和双峰自旋。

\emph{3第三步：Newman--Ruelle 零点聚积与相变。}Newman--Ruelle（1976）把 Lee--Yang 圆定理与热力学极限结合：若零点在单位圆上以密度 \(\rho(\theta)\) 聚积，则自由能 \(f(z)\) 在 \(|z|\ne1\) 处解析，在 \(|z|=1\) 上有奇性当且仅当 \(\rho(\theta)>0\)。对铁磁 Ising，\(\rho(0)>0\) 蕴含 \(z=1\)（\(h=0\)）处相变。更精确地，磁化
\begin{equation*}
  m=\pi\rho(0),
\end{equation*}
把序参量直接与零点密度联系。这一公式是 Yang 1952 年磁化公式的零点密度版本：磁化非零当且仅当零点在物理点聚积。

\emph{第四步：非铁磁耦合的失效与补救。}反铁磁耦合（\(J_{ij}<0\)）破坏圆定理：零点不再限于单位圆，Asano--Ruelle 收缩的分离条件失效。Biskup--Chayes--Crawford（2008）在特定条件下恢复反铁模零点控制。对 frustrated 系统（耦合符号混合），一般无零点定理，零点可分布在复平面任意位置——这是 spin glass 复杂性的体现。对 \(O(n)\) 模型（\(n\ge2\)），Lee--Yang 性质一般失效（连续对称性使零点结构更复杂），但 Michel--Sénéchal 给出了特定 \(O(n)\) 模型的-零点约束。这些推广和D失败把 Lee--Yang 理论的适用边界精确划定：铁磁 Ising 型是充分条件，不是必要条件。






















\emph{[一维 Ising 的精确解与转移矩阵谱]}
一维 Ising 模型虽无相变，但其精确解提供了转移矩阵方法的完整原型。分三步给出全部推导。

\emph{第一步：转移矩阵与精确配分函数。}对 \(N\) 自旋周期链（\(\sigma_{N+1}=\sigma_1\)），Hamiltonian \(H=-J\sum_{i=1}^N\sigma_i\sigma_{i+1}-h\sum_{i=1}^N\sigma_i\)。转移矩阵
\begin{equation*}
  \mathbf T=\begin{pmatrix}\ee^{K+H}&\ee^{-K}\\\ee^{-K}&\ee^{K-H}\end{pmatrix},\qquad K=\beta J,\;H=\beta h.
\end{equation*}
配分函数 \(Z_N=\Tr\mathbf T^N=\lambda_+^N+\lambda_-^N\)，本征值
\begin{equation*}
  \lambda_\pm=\ee^K\cosh H\pm\sqrt{\ee^{2K}\sinh^2 H+\ee^{-2K}}.
\end{equation*}
热力学极限 \(N\to\infty\)：\(-\beta f=\log\lambda_+\)（\(\lambda_+>\lambda_-\) 始终成立）。

\emph{第二步：自由能与磁化。}零场（\(H=0\)）：\(\lambda_\pm=2\cosh K\pm2\sinh K=\ee^{\pm K}\cdot2\cosh K\)，更精确地 \(\lambda_\pm=2\cosh K\pm2\sinh K\)，故 \(\lambda_+=2\cosh K+2\sinh K=2\ee^K\)，\(\lambda_-=2\ee^{-K}\)。自由能
\begin{equation*}
  -\beta f=\log(2\cosh K)+\log\!\left(1+\sqrt{1-\tanh^2 K}\right)=\log(2\cosh K)+\log\!\left(1+\frac{1}{\cosh K}\right).
\end{equation*}
更简洁地：\(-\beta f=\log\lambda_+=\log(2\cosh K)\)（零场下 \(\lambda_+=2\cosh K\)，因 \(\sqrt{\ee^{-2K}\cdot0+\ee^{-2K}}=\ee^{-K}\)，\(\lambda_+=\ee^K\cosh 0+\sqrt{\ee^{2K}\cdot0+\ee^{-2K}}=\ee^K+\ee^{-K}=2\cosh K\)）。磁化
\begin{equation*}
  m=-\frac{\partial f}{\partial h}=\frac{\sinh H}{\sqrt{\sinh^2 H+\ee^{-4K}}},
\end{equation*}
零场极限 \(H\to0\)：\(m\to\frac{H}{\sqrt{H^2+\ee^{-4K}}}\to0\)（\(K<\infty\)），无自发磁化。

\emph{第三步：关联函数与相关长度。}二点关联 \(\langle\sigma_0\sigma_r\rangle\) 由次主本征值控制：
\begin{equation*}
  \langle\sigma_0\sigma_r\rangle=\left(\frac{\lambda_-}{\lambda_+}\right)^r=(\tanh K)^r=\ee^{-r/\xi},
\end{equation*}
相关长度 \(\xi^{-1}=-\log\tanh K\)。\(K\to\infty\)（\(T\to0\)）：\(\tanh K\to1\)，\(\xi\to\infty\)（零温相变）。\(K\to0\)（\(T\to\infty\)）：\(\tanh K\to0\)，\(\xi\to0\)（完全无序）。有限温度下 \(\xi<\infty\)，关联指数衰减，无长程序。一维 Ising 的无相变性因此精确可证：\(\lambda_+/\lambda_-=\coth K>1\) 始终成立，谱隙不闭合，Perron--Frobenius 严格占优。










